% figures/logic_gates.tex
% Ręcznie rysowane bramki w stylu ANSI/US.
% Współrzędne: (x rośnie w prawo, y rośnie w górę).
\begin{tikzpicture}[
    thick, font=\small,
    >={Latex[length=2mm]}
]

% =========================================================
%  WEJŚCIA  S1..S4
% =========================================================
\node (S1) at (0, 4.4) {$S_1$};
\node (S2) at (0, 3.6) {$S_2$};
\node (S3) at (0, 1.4) {$S_3$};
\node (S4) at (0, 0.6) {$S_4$};

% =========================================================
%  BRAMKA OR (4-wejściowa)
%  Środek pionowy: y=2.5, lewa krawędź x=2.0, szer. 1.6, pół-wys. 1.4
% =========================================================
\def\orL{2.0}
\def\orC{2.5}
\def\orW{1.6}
\def\orH{1.4}

% Korpus bramki OR: górna krawędź, dolna krawędź i łuk z tyłu
\draw (\orL, \orC+\orH)
    .. controls (\orL+0.6, \orC+\orH) and (\orL+1.1, \orC+0.6)
    .. (\orL+\orW, \orC)
    .. controls (\orL+1.1, \orC-0.6) and (\orL+0.6, \orC-\orH)
    .. (\orL, \orC-\orH)
    .. controls (\orL+0.5, \orC-0.7) and (\orL+0.5, \orC+0.7)
    .. (\orL, \orC+\orH);

% Wprowadzenia S1..S4 do bramki OR
\draw (S1.east) -- (2.55, 4.4) -- (2.55, 3.5);
\draw (S2.east) -- (2.55, 3.6) -- (2.55, 3.0);
\draw (S3.east) -- (2.55, 1.4) -- (2.55, 2.0);
\draw (S4.east) -- (2.55, 0.6) -- (2.55, 1.5);

% Wyjście OR -> P
\draw (\orL+\orW, \orC) -- (5.2, \orC);
\node[above] at (4.5, \orC) {$P$};

% =========================================================
%  BRAMKA XOR (2-wejściowa)
%  Środek: (6.2, 2.5), szer. 1.2, pół-wys. 0.9
% =========================================================
\def\xorL{5.6}
\def\xorC{2.5}
\def\xorW{1.2}
\def\xorH{0.9}

% Korpus XOR (identyczny jak OR)
\draw (\xorL, \xorC+\xorH)
    .. controls (\xorL+0.45, \xorC+\xorH) and (\xorL+0.8, \xorC+0.4)
    .. (\xorL+\xorW, \xorC)
    .. controls (\xorL+0.8, \xorC-0.4) and (\xorL+0.45, \xorC-\xorH)
    .. (\xorL, \xorC-\xorH)
    .. controls (\xorL+0.4, \xorC-0.45) and (\xorL+0.4, \xorC+0.45)
    .. (\xorL, \xorC+\xorH);

% Dodatkowy łuk z przodu (charakterystyczny dla XOR)
\draw (\xorL-0.22, \xorC+\xorH)
    .. controls (\xorL+0.18, \xorC+0.4) and (\xorL+0.18, \xorC-0.4)
    .. (\xorL-0.22, \xorC-\xorH);

% Wejście P do XOR
\draw (5.2, \orC) -- (\xorL-0.22, \orC) -- (\xorL-0.22, \xorC+0.4) -- (\xorL+0.3, \xorC+0.4);

% Wejście Q do XOR
\node (Qin) at (5.6, 4.6) {$Q$};
\draw (Qin.south) -- (5.6, 3.8) -- (\xorL-0.22, 3.8) -- (\xorL-0.22, \xorC-0.4) -- (\xorL+0.3, \xorC-0.4);

% Wyjście XOR
\draw (\xorL+\xorW, \xorC) -- (8.3, \xorC);

% =========================================================
%  BRAMKA AND (2-wejściowa)
%  Środek: (9.4, 2.5), szer. 1.2, pół-wys. 0.9
% =========================================================
\def\andL{8.3}
\def\andC{2.5}
\def\andW{1.2}
\def\andH{0.9}

% Korpus AND: prostokąt + półkole z prawej strony
\draw (\andL, \andC+\andH) -- (\andL+0.6, \andC+\andH)
      arc[start angle=90, end angle=-90, radius=0.9];
\draw (\andL, \andC-\andH) -- (\andL+0.6, \andC-\andH);

% Wejście 1 (z XOR)
\draw (8.3, \xorC) -- (\andL+0.3, \andC+0.4);
% (poprawka: wyrównujemy do poziomu wejścia AND; XOR wyjście jest na y=2.5,
%  więc trzeba lekko zejść)
\draw[white, line width=3pt] (8.3, \xorC) -- (\andL+0.3, \andC+0.4); % czyścimy poprzednią
\draw (8.3, \xorC) -- (8.5, \xorC) -- (8.5, \andC+0.4) -- (\andL+0.3, \andC+0.4);

% Wejście 2 (z NOT T)
\draw (\andL+0.3, \andC-0.4) -- (7.8, \andC-0.4) -- (7.8, 0.9);

% =========================================================
%  BRAMKA NOT (dla T) - trójkąt z kółkiem
% =========================================================
\def\notX{6.9}
\def\notY{0.9}
\draw (\notX, \notY+0.35) -- (\notX, \notY-0.35) -- (\notX+0.7, \notY) -- cycle;
\draw (\notX+0.7, \notY) circle[radius=0.08];

% Wejście T
\node (T) at (\notX-0.5, \notY-1.2) {$T$};
\draw (T.north) -- (\notX, \notY-1.2) -- (\notX, \notY);

% Wyjście NOT -> AND wejście 2
\draw (\notX+0.78, \notY) -- (7.8, \notY) -- (7.8, \andC-0.4) -- (\andL+0.3, \andC-0.4);

% =========================================================
%  WYJŚCIE
% =========================================================
\draw (\andL+0.6+0.9, \andC) -- (10.8, \andC) node[right] {$Q_{\text{next}}$};

\end{tikzpicture}